problem:  
Let \((X,d,\mu)\) be a proper geodesic metric measure space satisfying the **weak Bishop–Gromov volume monotonicity**: there exist parameters \(k\in \mathbb{R}\) and \(N>0\) such that, for every \(x\in X,\) the function  
\[
r\mapsto \frac{\mu(B(x,r))}{v_{k,N}(r)}
\]
is nonincreasing for all sufficiently small \(r>0,\) where \(v_{k,N}(r)\) is the volume of a radius‑\(r\) ball in the \(N\)-dimensional model space of constant sectional curvature \(k\). Suppose also that for \(\mu\)-almost every \(x\in X\), there exists at least one tangent cone \((Y_x, d_x, \nu_x, y_x)\) in the pointed measured Gromov–Hausdorff sense.

Open problem:  
Does the weak Bishop–Gromov monotonicity, without any auxiliary curvature-dimension (\(CD(K,N)\)) or infinitesimal Hilbertian assumptions, guarantee **uniqueness of the tangent measure class** at \(\mu\)-almost every \(x\in X\)?  

Equivalently: must the set of all tangent measures \(\mathsf{Tan}(\mu,x)\) at almost every \(x\) form a *single homothetic class*—that is, is there a constant \(c_x>0\) and a measure \(\nu_x\) such that every tangent measure at \(x\) equals \(c_x\,(\mathrm{dil}_r)_\#\nu_x\) for some rescaling factor \(r>0\)?  
If not, can one construct a counterexample metric measure space that satisfies the Bishop–Gromov monotonicity yet admits nonunique tangent measures at positive \(\mu\)-measure of points?

Why is it a "good" problem:  
This problem examines whether **Bishop–Gromov monotonicity alone**, which is a global comparison inequality, already enforces a *local uniqueness property*—the determinacy of infinitesimal measure geometry. It strips away the heavy analytical machinery of curvature‑dimension and RCD theories to test the conceptual limit of what volume comparison alone can dictate.

- **Profound insight:** It directly addresses whether *macroscopic curvature control through volume growth* determines the *microscopic geometric structure* of the measure. If true, it would identify Bishop–Gromov monotonicity as a cornerstone principle shaping infinitesimal regularity beyond smooth or analytic contexts.  
- **Bridge between fields:** It links comparison geometry, tangent‑measure theory (à la Preiss), and analysis on metric spaces, probing whether monotone volume ratios can substitute for PDE-based curvature constraints in producing rigidity of tangent cones.  
- **Extreme simplicity:** The statement—asking if Bishop–Gromov monotonicity forces uniqueness of tangent measures—uses only fundamental geometric notions: balls, volumes, rescaling, and limits. Yet its resolution would require novel interaction of measure theory, geometry, and analysis at the singular boundary of curvature comparison.